% Original: PDF strana 9, gornji slajd.
\slajdid{02-s017}
\begin{frame}[label=02-s017]{Kanonička konjunktivna forma}
% element: 02-s017:e01
\[F=(C+B+A)(C+B+\bar A)(C+\bar B+A)(\bar C+B+A)(\bar C+\bar B+\bar A).\]
\par\vspace{3mm}
% element: 02-s017:e02
\par\Rezultat{Proizvod potpunih zbirova naziva se \alert{normalna, odnosno kanonička konjunktivna forma}.}
\par\vspace{4mm}
% element: 02-s017:e03
\Oznaka{Pravilo iz reda sa \(F=0\)}
Ulaz 0 daje promenljivu bez crte; ulaz 1 daje komplementnu promenljivu. Tako dobijeni zbir je nula u tom redu. Svi zbirovi povezuju se operacijom I.
\end{frame}
